\documentclass[11pt,letterpaper]{article}
\usepackage[margin=1in]{geometry}
\usepackage{graphicx} % Required for inserting images
\usepackage{amsthm,amsmath,amssymb}
\usepackage{algorithm}
\usepackage{algpseudocode}
\usepackage{xcolor}
\usepackage{soul}
\usepackage{enumitem}

\usepackage[hidelinks]{hyperref}
\usepackage{cleveref}

% Theorems
\newtheorem{theorem}{Theorem}[section]
\newtheorem{corollary}[theorem]{Corollary}
\newtheorem{lemma}[theorem]{Lemma}
\newtheorem{claim}[theorem]{Claim}
\theoremstyle{definition}
\newtheorem{definition}[theorem]{Definition}
\newtheorem{remark}[theorem]{Remark}
\newtheorem{observation}[theorem]{Observation}

\crefname{theorem}{Theorem}{Theorems}
\Crefname{theorem}{Theorem}{Theorems}
\crefname{corollary}{Corollary}{Corollaries}
\Crefname{corollary}{Corollary}{Corollaries}
\crefname{lemma}{Lemma}{Lemmas}
\Crefname{lemma}{Lemma}{Lemmas}
\crefname{claim}{Claim}{Claims}
\Crefname{claim}{Claim}{Claims}
\crefname{definition}{Definition}{Definitions}
\Crefname{definition}{Definition}{Definitions}
\crefname{remark}{Remark}{Remarks}
\Crefname{remark}{Remark}{Remarks}
\crefname{observation}{Observation}{Observations}
\Crefname{observation}{Observation}{Observations}

\newcommand{\reals}{\mathbb{R}}
\newcommand{\Span}{\mathrm{span}}
\newcommand{\coreq}{\mathrm{coreq}}

\newcommand{\calA}{\mathcal{A}}
\newcommand{\calB}{\mathcal{B}}
\newcommand{\calC}{\mathcal{C}}
\newcommand{\calD}{\mathcal{D}}
\newcommand{\calE}{\mathcal{E}}
\newcommand{\calF}{\mathcal{F}}
\newcommand{\calG}{\mathcal{G}}
\newcommand{\calH}{\mathcal{H}}
\newcommand{\calI}{\mathcal{I}}
\newcommand{\calJ}{\mathcal{J}}
\newcommand{\calK}{\mathcal{K}}
\newcommand{\calL}{\mathcal{L}}
\newcommand{\calM}{\mathcal{M}}
\newcommand{\calN}{\mathcal{N}}
\newcommand{\calO}{\mathcal{O}}
\newcommand{\calP}{\mathcal{P}}
\newcommand{\calQ}{\mathcal{Q}}
\newcommand{\calR}{\mathcal{R}}
\newcommand{\calS}{\mathcal{S}}
\newcommand{\calT}{\mathcal{T}}
\newcommand{\calU}{\mathcal{U}}
\newcommand{\calV}{\mathcal{V}}
\newcommand{\calW}{\mathcal{W}}
\newcommand{\calX}{\mathcal{X}}
\newcommand{\calY}{\mathcal{Y}}
\newcommand{\calZ}{\mathcal{Z}}

\newcommand{\cost}{\mathrm{COST}}
\newcommand{\ALG}{\mathrm{ALG}}
\newcommand{\OPT}{\mathrm{OPT}}
\newcommand{\Flow}{\textsc{Flow}}
\newcommand{\ouralg}{\textsc{LineTracker}}

\newcommand{\EE}{\mathbb{E}}
\newcommand{\Conv}{\mathrm{Conv}}

\newcommand{\interval}[1]{\lfloor #1 \rceil}
\newcommand{\dis}[1]{#1}
\newcommand{\off}{\mathrm{off}}
\newcommand{\on}{\mathrm{on}}
\newcommand{\org}{\mathrm{org}}
\newcommand{\rel}{\mathrm{rel}}
\newcommand{\new}{\mathrm{new}}

\newcommand{\solution}{\medskip\noindent\textbf{Solution.}\par\medskip}

\title{47-834: Linear Programming \\ Homework 1}
\author{Poon Tze-Yang}
\date{Due September 15, 2026}


\begin{document}

\maketitle

\noindent Fall 2026 Mini-1

\begin{enumerate}
\item \textbf{[15 points]} Consider the problem
\begin{align*}
    \text{minimize} \quad & \frac{c^\top x + d}{f^\top x + g} \\
    \text{subject to} \quad & Ax \leq b \\
    & f^\top x + g > 0.
\end{align*}
\begin{enumerate}[label=(\roman*)]
\item Suppose that we have some prior knowledge that the optimal cost belongs to an interval $[K,L]$. Provide a procedure that allows us to compute the optimal cost within any desired accuracy.
\item Suppose that this optimization problem is solvable. Show a way to obtain a lower bound $K$ and an upper bound $L$ for the optimal cost.
\end{enumerate}
For both of the questions above, you can use linear programming as a subroutine.
\end{enumerate}

\solution
\paragraph{(i)} We can perform binary search on the interval.
Concretely, given an interval $[K,L]$, let its midpoint be $M$.
If the minimum value is less than or equal to $M$, then the following program is feasible:
\begin{align*}
    \text{find} \quad & x \\
    \text{subject to} \quad & Ax \leq b \\
    & c^\top x + d - M(f^\top x + g) \leq 0 \\
    & f^\top x + g > 0,
\end{align*}
and we can update the interval to $[K,M]$.
Otherwise, we update the interval to $[M,L]$.
We can repeat this process until the length of the interval is less than the desired accuracy.

\paragraph{(ii)} Any feasible solution $x$ gives an upper bound $L$. 
We can obtain one by solving for LP feasibility without the objective value.

We can obtain a lower bound by separately minimising the numerator and maximizing the denominator. 
Concretely, we can solve the following two linear programs:
\begin{align*}
    \text{minimize} \quad & c^\top x + d \\
    \text{subject to} \quad & Ax \leq b \\
    & f^\top x + g > 0,
\end{align*}
and
\begin{align*}
    \text{maximize} \quad & f^\top x + g \\
    \text{subject to} \quad & Ax \leq b \\
    & f^\top x + g > 0. 
\end{align*}
Let the optimal values of these two programs be $N$ and $D$ respectively.
Then we can take $K = N/D$ as a lower bound for the optimal cost.

\begin{enumerate}[resume]
\item \textbf{[15 points]} Consider a polyhedron $P$ given in inequality form, i.e., $P = \{x\in\reals^n : a_i^\top x \leq b_i \ \forall i = 1,\ldots,m\}$. A Euclidean ball with center $y\in\reals^n$ and radius $r$ is defined as $B_r(y) := \{x\in\reals^n : \|x-y\|_2 \leq r\}$, i.e., set of all points within Euclidean distance of $r$ from $y$. We are interested in finding a ball with largest possible radius, which is entirely contained in $P$. The center of such a ball is called Chebychev center of $P$. Provide a linear programming formulation of this problem.
\end{enumerate}

\solution
First, we observe that given the center $y$ of such a ball of radius $r$, tightening all constraints of $P$ by $r$ preserves the feasibility of $c$.
Hence, we can try to maximize $r$ subject to the feasibility constraints.
This can be formulated as follows:
\begin{align*}
    \text{maximize} \quad & r \\
    \text{subject to} \quad & a_i^\top y + r \leq b_i \quad \forall i = 1,\ldots,m \\
    & r \geq 0.
\end{align*}    
The maximizing $y$ gives the Chebychev center of $P$.

\begin{enumerate}[resume]
\item \textbf{[15 points]} Given two sets $U, V\subset\reals^m$, we define $U + V = \{x\in\reals^m : \exists u\in U, \exists v\in V \text{ s.t. } x = u+v\}$.

Let $D := \{x\in\reals^n : Ax + b + Q_s \subseteq P \ \forall s\in S\}$ where the set $\emptyset \neq P \subset \reals^m$ admits polyhedral representation, the set $\emptyset\neq S\subset\reals^k$ is given but arbitrary, and the sets $\emptyset \neq Q_s \subset \reals^m$ are indexed by $s\in S$.
\begin{enumerate}
\item Suppose that $S$ is a finite set and for each $s\in S$ we have $Q_s = \{q_s\}$, i.e., is a single point. Then, will the set $D$ be polyhedrally representable?
\item State sufficient conditions on the structure of sets $Q_s$ and $S$ that will guarantee that the resulting set $D$ is polyhedral. Here, the goal is to have conditions as general as possible. Among your sufficient conditions, can you identify at least some of those that are necessary?
\end{enumerate}
\end{enumerate}

\solution

\paragraph{(a)} Yes, the set $D$ will be polyhedrally representable.
For each $s\in S$, we have $Ax + b + q_s \subseteq P$. 
This defines a polyhedron, and $D$ is the intersection of these polyhedra.
Hence, $D$ is also a polyhedron.

\paragraph{(b)} In fact, no conditions on $S$ or $Q_s$ are necessary to guarantee that $D$ is polyhedral.
First, note that the affine transformation $Ax + b$ preserves polyhedrality.
Furthermore, taking only one $Q$ which is the union of all $Q_s$ is equivalent.
Thus, $D$ is polyhedral if and only of $D' = \{x\in\reals^n : x + Q \subseteq P\}$ is polyhedral.
For each constraint of $P$, the set of values of $x$ that satisfy $x + Q \subseteq P$ forms a half-space, since only the largest $Q$ in the direction normal to that constraint matters.
Hence, $D'$ is the intersection of half-spaces, and is polyhedral.

\begin{enumerate}[resume]
\item \textbf{[15 points]} Which of the following claims is always true? Give a proof of the correct claim and give a counter example for incorrect one.
\begin{enumerate}
\item The set $X = \{x\in\reals^n : Ax\leq b\}$ is a cone if and only if $X = \{x\in\reals^n : Ax\leq 0\}$.
\item The set $X = \{x\in\reals^n : Ax\leq b\}$ is a cone if and only if $b = 0$.
\end{enumerate}
\end{enumerate}

\solution

\paragraph{(a)} The claim is true.
If the set $X$ is a cone, then the zero vector must be in $X$, which implies that $b\geq 0$.
Furthermore, if $b_i>0$, then for the corresponding row $a_i$ of $A$, $a_i$ must be zero, otherwise we have a contradiction: $a_i$ is in $X$ and since $X$ is a cone, we must be able to scale $a_i$ up until we violate the constraint $a_i^\top x \leq b$.
Thus, for each row, we either have that $b_i = 0$ or $a_i = 0$ and setting $b_i$ to zero does not change the set $X$.

\paragraph{(b)} The claim is not true. 
We can add redundant constraints to the system $Ax\leq 0$ to obtain a system $A'x\leq b'$ with $b'\neq 0$ that defines the same set.

\begin{enumerate}[resume]
\item \textbf{[20 points]} Prove the following statements:
\begin{enumerate}
\item If a system of $m$ linear equations $Ax = b$ with $n$ variables (i.e., $x\in\reals^n$) has a nonnegative solution, it has a nonnegative solution with at most $m$ positive entries.
\item If a system of $m$ linear inequalities $Ax\leq b$ with $n$ variables (i.e., $x\in\reals^n$) is infeasible, so is a properly selected sub-system composed of at most $n+1$ inequalities from the system.
\end{enumerate}
\end{enumerate}

\solution
\paragraph{(a)} Considering the system of equations with the non-negativity constraints, we are guaranteed to have feasible solutions.
At basic feasible solutions, at least $n-m$ of the non-negativity constraints are tight.
Hence, at least $n-m$ of the variables are zero, and at most $m$ of the variables are positive.

\paragraph{(b)} We prove the contrapositive, which states that if every sub-system of $n+1$ inequalities from the system is feasible, then the original system is feasible.
Note that the half-space defined by each inequality is convex.
Hence, by Helly's theorem, if every sub-system of $n+1$ inequalities is feasible, then the intersection of all half-spaces is non-empty, and the original system is feasible.


\begin{enumerate}[resume]
\item \textbf{[20 points]} Let $V_1,\ldots,V_n$ be $n$ nonempty sets in $\reals^m$, and define
\[
V := \Conv(V_1 + V_2 + \ldots + V_n).
\]
Prove the following statements:
\begin{enumerate}
\item $V = \Conv(V_1) + \ldots + \Conv(V_n)$.
\item For any $x\in V$, there exists a representation of $x$ such that
\[
x = x_1 + \ldots + x_n, \qquad \text{where } x_i \in \Conv(V_i),
\]
in which at least $n-m$ of $x_i$'s belong to the respective sets $V_i$.

Informally, this result states that when $n\gg m$, summing up $n$ nonempty sets in $\reals^m$ possesses certain ``convexification property'' -- every point from the convex hull $V$ of the sum of our sets is the sum of points $x_i$ with all but $m$ of them belonging to $V_i$ rather than to $\Conv(V_i)$, and only $\leq m$ of the points belonging to $V_i$ ``fractionally,'' that is, belonging to $\Conv(V_i)$, but not to $V_i$. This nice fact has a lot of useful applications.
\end{enumerate}
\end{enumerate}

\solution

\paragraph{(a)}

\ul{$V\subseteq \Conv(V_1) + \ldots + \Conv(V_n)$:} 
For all $v\in V$, we can express $v = \sum_{j=1}^\ell \lambda_j v_j$ as a convex combination of elements of $V_1 + \dots + V_n$.
For each $v_j$, we can express $v_j = \sum_{i=1}^n v_{j,i}$, where $v_{j,i}\in V_i$.
Thus, we have $v = \sum_{j=1}^\ell \lambda_j \sum_{i=1}^n v_{j,i} = \sum_{i=1}^n \sum_{j=1}^\ell \lambda_j v_{j,i}$.
For each $i$, we have $\sum_{j=1}^\ell \lambda_j v_{j,i} \in \Conv(V_i)$, and hence $v\in \Conv(V_1) + \ldots + \Conv(V_n)$.

\ul{$\Conv(V_1) + \ldots + \Conv(V_n) \subseteq V$:} 
For all $v\in \Conv(V_1) + \ldots + \Conv(V_n)$, we can express $v = \sum_{i=1}^n v_i$, where $v_i\in \Conv(V_i)$.
For each $v_i$, we can express $v_i = \sum_{j=1}^{\ell_i} \lambda_{i,j} v_{i,j}$ as a convex combination of elements of $V_i$.
Thus, we have $v = \sum_{i=1}^n \sum_{j=1}^{\ell_i} \lambda_{i,j} v_{i,j}$.
This can be converted to a convex combination of elements of $V_1 + \dots + V_n$ by repeatedly taking elements of $V_1 + \dots + V_n$ corresponding to the smallest current $\lambda_{i,j}$.

\paragraph{(b)}
Using (a), for each $x\in V$, we can express $x = \sum_{i=1}^n x_i$, where $x_i\in \Conv(V_i)$.
Using Caratheodory's Theorem, each $x_i$ can be expressed as the convex combination of finitely many ($m+1$) points in $V_i$:
\begin{align*}
    x_i = \sum_{j\in [m+1]}\lambda_{i,j} x_{i,j}.
\end{align*}
Treating the points $x_{i,j}$ as fixed and $\lambda_{i,j}$ as the variables in $\geq 0$, we can rewrite the contraints on $\lambda$ as follows:
\begin{align*}
    &\sum_{i\in[n]}\sum_{j\in [m+1]} \lambda_{i,j}x_{i,j} = x\\
    &\sum_{j\in[m+1]}\lambda_{i,j} = 1 &\forall i\in [n].
\end{align*}
The first set of equations is $m$ equations because $x$ has $m$ dimensions, and the second set of equations is $n$ equations.
Let the number of $\lambda$s be $\ell$. 
Then an extreme point, we need at least $\ell - n - m$ of the $\lambda$'s to be zero (tight $\lambda_{i,j}\geq 0$ constraint).
Thus, at most $n+m$ of the $\lambda$s are positive.
Furthermore, we need at least one positive lambda for each of the $n$ sets $V_i$, and if one of the $\lambda_{i,j}$ for $V_i$ is fractional, at least two $\lambda_{i,j}$'s for that $V_i$ must be fractional.
Hence, at most $m$ of the $V_i$'s can have fractional $\lambda_{i,j}$'s, so at least $n-m$ of the $V_i$'s have only one positive $\lambda_{i,j}$ with value equal to 1.
This gives the claim.



\vspace{\baselineskip}

The following question is extra credit!

\begin{enumerate}[resume]
\item \textbf{(Extra Credit) [20 points]} An asteroid field contains 400 tons of a rare mineral. You may model each mineral-bearing asteroid as a small three-dimensional ball of positive radius, with the mineral distributed throughout the asteroid.

Two rival space agencies, Astra and Borealis, agree to divide the mineral resources according to the following rather unusual treaty.
\begin{itemize}
\item First, Astra chooses a point $x\in\reals^3$ at which to place a space station.
\item Then Borealis chooses an 2D arbitrary plane $\Pi$ passing through $x$, together with one of the two sides of $\Pi$. Borealis receives all of the mineral lying on its chosen side, while Astra receives all of the mineral on the other side.
\end{itemize}
\begin{enumerate}
\item \textbf{A bad asteroid field:} Construct an arrangement of the asteroids for which Astra cannot guarantee receiving more than 101 tons of of mineral, regardless of where it places its station.
\item \textbf{A guaranteed safe location:} Astra can always choose the location $x$ so that, no matter which plane and side Borealis chooses, Astra receives at least 99 tons of mineral.
\item \textbf{An $n$-dimensional universe:} Now imagine that the asteroid field lies in $\reals^n$, contains a mineral mass of $M$, and Astra again chooses a point $x$. Borealis then chooses a nonzero vector $e$ and one of the two half-spaces determined by
\[
e^\top (x-a) = 0.
\]
Prove that for every $\epsilon > 0$, Astra can place its station so as to guarantee at least $\frac{M}{n+1} - \epsilon$ units of mineral. Also, show that, in general, Astra cannot guarantee $\frac{M}{n+1} + \epsilon$ units of mineral.
\end{enumerate}
If you submit an answer to this question, I will ask you to present your answer in person!
\end{enumerate}

\solution

\paragraph{(a)}
Consider the following asteroid field: place one clump of 100 minerals at the origin, and 100 minerals at a distance $d>10$ from the origin in the positive direction on each axis.
Assume without loss of generality that each clump has radius 1.
If Astra chooses a point for the space station that is greater than 1 in any coordinate, say, $x$ , Borealis can choose the hyperplane $x=1$ and choose the side with the origin for at least 300 tons.
If Astra chooses a point with each coordinate smaller than $2$, Borealis can choose the plane $x+y+z= r$ for the appropriate $r$ and choose the side not containing the origin for at least 300 tons.
This covers all points Astra can choose.

\paragraph{(b)}
(This seems to need Helly's theorem with infinite sets, which requires compactness)
% We will use Helly's theorem.
% Consider the family of subsets of asteroids that can be contained in half-spaces which fully contains at least 301 mineral-bearing asteroids.
% There are finitely many such subsets of asteroids since there are finitely many asteroids.
% Also note that any 4 such subsets intersect at at least one asteroid.
% Hence, by Helly's theorem, there is at least one asteroid that is containted in every such subset.

% Astra can place the space station at this asteroid.
% Consider any hyperplane passing through this asteroid.


\paragraph{(c)}

(The $\varepsilon$ deals with the points along the plane.)
A similar construction as in (a) works to show that $\frac{M}{n+1} + \varepsilon$ units is not achievable, where we place one cluster of units at the origin and one cluster along each coordinate.


\end{document}
